\documentclass{article}
\usepackage{graphicx} % Required for inserting images
\usepackage{amsmath}
\usepackage{amssymb}
\title{Lesson 26 Homework - Proof Writing}
\author{William Wang}
\date{March 2026}

\begin{document}

\maketitle

\section{}
\subsection{}
\[\begin{pmatrix}
    3&1&-2&|&-7 \\ 4&2&2&|&8 \\ -1&-1&3&|&6
\end{pmatrix}\]
Then \
\(R_1 \rightarrow \frac{1}{3}R_1\)
\[\begin{pmatrix}
    1&1/3&-2/3&|&-7/3 \\ 4&2&2&|&8 \\ -1&-1&3&|&6
\end{pmatrix}\]
Then \(R_2 \rightarrow  R_2 - 4R_1\)
\[\begin{pmatrix}
    1&1/3&-2/3&|&-7/3 \\ 0&2/3&14/3&|&52/3 \\ -1&-1&3&|&6
\end{pmatrix}\]
Then \(R_3 \rightarrow R_3 + R_1\)
\[\begin{pmatrix}
    1&1/3&-2/3&|&-7/3 \\ 0&2/3&14/3&|&52/3 \\ 0&-2/3&7/3&|&11/3
\end{pmatrix}\]
Then \(R_2 \rightarrow \frac{2}{3}R_2\)
\[\begin{pmatrix}
    1&1/3&-2/3&|&-7/3 \\ 0&1&7&|&26 \\ 0&-2/3&7/3&|&11/3
\end{pmatrix}\]
Then \(R_3 \rightarrow R_3 + \frac{2}{3}R_2\)
\[\begin{pmatrix}
    1&1/3&-2/3&|&-7/3 \\ 0&1&7&|&26 \\ 0&0&7&|&21
\end{pmatrix}\]
Then \(R_3 \rightarrow \frac{1}{3}R_3\)
\[\begin{pmatrix}
    1&1/3&-2/3&|&-7/3 \\ 0&1&7&|&26 \\ 0&0&1&|&3
\end{pmatrix}\]
Then we have \(z = 3\). 
\[y+7z = 26 \implies y = 5\]
\[x + y/3 -2z/3 = -7/3 \implies x = -2\]
Hence the solution set is a single point at \((-2,5,3)\). 
\subsection{}
\[\begin{pmatrix}
    3&1&-2&|&2 \\ 1&-2&1&|&3 \\ -1&-1&5&|&4
\end{pmatrix}\]
Then \(R_1 \rightarrow \frac{1}{3}R_1\)
\[\begin{pmatrix}
    1&1/3&-2/3&|&2/3 \\ 1&-2&1&|&3 \\ -1&-1&5&|&4
\end{pmatrix}\]
Then \(R_2 \rightarrow R_2-R_1\)
\[\begin{pmatrix}
    1&1/3&-2/3&|&2/3 \\ 0&-7/3&5/3&|&7/3 \\ -1&-1&5&|&4
\end{pmatrix}\]
Then \(R_3 \rightarrow R_3 + R_1\)
\[\begin{pmatrix}
    1&1/3&-2/3&|&2/3 \\ 0&-7/3&5/3&|&7/3 \\ 0&-2/3&13/3&|&14/3
\end{pmatrix}\]
Then \(R_2 \rightarrow -\frac{3}{7}R_2\)
\[\begin{pmatrix}
    1&1/3&-2/3&|&2/3 \\ 0&1&-5/7&|&-1 \\ 0&-2/3&13/3&|&14/3
\end{pmatrix}\]
Then \(R_3 \rightarrow \frac{3}{2}R_3 + R_2\)
\[\begin{pmatrix}
    1&1/3&-2/3&|&2/3 \\ 0&1&-5/7&|&-1 \\ 0&0&81/14&|&6
\end{pmatrix}\]
Then \(R_3 \rightarrow \frac{14}{81}R_3\)
\[\begin{pmatrix}
    1&1/3&-2/3&|&2/3 \\ 0&1&-5/7&|&-1 \\ 0&0&1&|&28/27
\end{pmatrix}\]
Hence \(z = 28/27\). 
\[y-5z/7 = -1 \implies y = -7/27\]
\[x+y/3-2z/3 = 2/3 \implies x = 13/9\]
Hence the solution set is a single point at \((13/9, -7/27, 28/27)\)
\subsection{}
\[\begin{pmatrix}
    3&1&-2&|&2 \\ 1&-2&1&|&3 \\ -1&-5&4&|&46
\end{pmatrix}\]
Then \(R_1 \rightarrow \frac{1}{3}R_1\)
\[\begin{pmatrix}
    1&1/3&-2/3&|&2/3 \\ 1&-2&1&|&3 \\ -1&-5&4&|&46
\end{pmatrix}\]
Then \(R_2 \rightarrow R_2 - R_1\)
\[\begin{pmatrix}
    1&1/3&-2/3&|&2/3 \\ 0&-7/3&5/3&|&7/3 \\ -1&-5&4&|&46
\end{pmatrix}\]
Then \(R_3 \rightarrow R_3+R_1\)
\[\begin{pmatrix}
    1&1/3&-2/3&|&2/3 \\ 0&-7/3&5/3&|&7/3 \\ 0&-14/3&10/3&|&140/3
\end{pmatrix}\]
Then \(R_2 \rightarrow -\frac{3}{7}R_2\)
\[\begin{pmatrix}
    1&1/3&-2/3&|&2/3 \\ 0&1&-5/7&|&-1 \\ 0&-14/3&10/3&|&140/3
\end{pmatrix}\]
Then \(R_3 \rightarrow R_2 + \frac{3}{14}R_3\)
\[\begin{pmatrix}
    1&1/3&-2/3&|&2/3 \\ 0&1&-5/7&|&-1 \\ 0&0&0&|&9
\end{pmatrix}\]
The systems is inconsistent because of the last row. Hence there are no solutions. 
\section{}
\[\begin{cases}
    x+y+z = 1 \\ 2x+2y+2z = 2 \\ 3x+3y+3z=3
\end{cases}\]
The exact solution is a plane in \(R^3\)defined by the equation \(x+y+z = 1\). 
\section{}
\subsection{}
\[\begin{pmatrix}
    1&2&-3&|&u \\ 4&-2&16&|&v \\ -1&3&5&|&w
\end{pmatrix}\]
Then \(R_2 \rightarrow \frac{1}{4}R_2 - R_1\)
\[\begin{pmatrix}
    1&2&-3&|&u \\ 0&-5/2&-1&|&v/4 - u \\ -1&3&5&|&w
\end{pmatrix}\]
Then \(R_3 \rightarrow R_3 + R_1\)
\[\begin{pmatrix}
    1&2&-3&|&u \\ 0&-5/2&-1&|&v/4 - u \\ 0&5&2&|&u+w
\end{pmatrix}\]
Then \(R_2\rightarrow4R_2\)
\[\begin{pmatrix}
    1&2&-3&|&u \\ 0&-10&-4&|&v-4u \\ 0&5&2&|&u+w
\end{pmatrix}\]
Then \(R_3 \rightarrow 2R_3 +R_2\)
\[\begin{pmatrix}
    1&2&-3&|&u \\ 0&-10&-4&|&v-4u \\ 0&0&0&|&(v-4u) + 2(u+w)
\end{pmatrix}\]
For the system to be consistent, the last row cannot be \([0,0,0|\text{nonzero}]\). Hence we must have 
\[v-4u+2u+2w = 0\]
\[v-2u+2w = 0\]
\subsection{}
\[\begin{pmatrix}
    1&2&-3&|&u \\ 2&-1&-7&|&v \\ -1&3&5&|&w
\end{pmatrix}\]
Then \(R_2 \rightarrow \frac{1}{2}R_2 - R_1\)
\[\begin{pmatrix}
    1&2&-3&|&u \\ 0&-5/2&-1/2&|&v/2-u \\ -1&3&5&|&w
\end{pmatrix}\]
Then \(R_3 \rightarrow R_1+R_3\)
\[\begin{pmatrix}
    1&2&-3&|&u \\ 0&-5/2&-1/2&|&v/2-u \\ 0&5&2&|&u+w
\end{pmatrix}\]
Then \(R_2 \rightarrow 2R_2\)
\[\begin{pmatrix}
    1&2&-3&|&u \\ 0&-5&-1&|&v-2u \\ 0&5&2&|&u+w
\end{pmatrix}\]
Then \(R_3 \rightarrow R_3 + R_2\)
\[\begin{pmatrix}
    1&2&-3&|&u \\ 0&-5&-1&|&v-2u \\ 0&0&1&|&v-u+w
\end{pmatrix}\]
Since the there is no row such that \([0,0,0|\text{nonzero}]\), the system is consistent for all values of \(v-u+w\). 
\section{}
\subsection{}
\[\begin{pmatrix}
    1&0&-3&|&1 \\ 1&2&-12&|&0 \\ 1&-1&2&|&0
\end{pmatrix}\]
Then \(R_2 \rightarrow R_2 - R_1\)
\[\begin{pmatrix}
    1&0&-3&|&1 \\ 0&2&-9&|&-1 \\ 1&-1&2&|&0
\end{pmatrix}\]
Then \(R_3 \rightarrow R_3 - R_1\)
\[\begin{pmatrix}
    1&0&-3&|&1 \\ 0&2&-9&|&-1 \\ 0&-1&5&|&-1
\end{pmatrix}\]
Then \(R_2 \rightarrow \frac{1}{2}R_2\)
\[\begin{pmatrix}
    1&0&-3&|&1 \\ 0&1&-9/2&|&-1/2 \\ 0&-1&5&|&-1
\end{pmatrix}\]
Then \(R_3 \rightarrow R_2 + R_3\)
\[\begin{pmatrix}
    1&0&-3&|&1 \\ 0&1&-9/2&|&-1/2 \\ 0&0&1/2&|&-3/2
\end{pmatrix}\]
Then \(R_3\rightarrow2R_3\)
\[\begin{pmatrix}
    1&0&-3&|&1 \\ 0&1&-9/2&|&-1/2 \\ 0&0&1&|&-3
\end{pmatrix}\]
Then we have \(c = -3\). 
\[b-9c/2=-1/2 \implies b = -14\]
\[a-3c = 1 \implies a= -8\]
\[(a,b,c)= (-8,-14,-3)\]
\subsection{}
\[\begin{pmatrix}
    1&0&-3&|&0 \\ 1&2&-12&|&1 \\ 1&-1&2&|&0
\end{pmatrix}\]
Then \(R_2 \rightarrow R_2 - R_1\)
\[\begin{pmatrix}
    1&0&-3&|&0 \\ 0&2&-9&|&1 \\ 1&-1&2&|&0
\end{pmatrix}\]
Then \(R_3 \rightarrow R_3 - R_1\)
\[\begin{pmatrix}
    1&0&-3&|&0 \\ 0&2&-9&|&1 \\ 0&-1&5&|&0
\end{pmatrix}\]
Then \(R_3 \rightarrow 2R_3 + R_2\)
\[\begin{pmatrix}
    1&0&-3&|&0 \\ 0&2&-9&|&1 \\ 0&0&1&|&1
\end{pmatrix}\]
Hence \(f = 1\). 
\[2e-9f = 1 \implies e = 5\]
\[d-3f = 0 \implies d = 3\]
\[(d,e,f) = (3,5,1)\]
\subsection{}
\[\begin{pmatrix}
    1&0&-3&|&0 \\ 1&2&-12&|&0 \\ 1&-1&2&|&1
\end{pmatrix}\]
Then \(R_2 \rightarrow R_2 - R_1\)
\[\begin{pmatrix}
    1&0&-3&|&0 \\ 0&2&-9&|&0 \\ 1&-1&2&|&1
\end{pmatrix}\]
Then \(R_3 \rightarrow R_3 - R_1\)
\[\begin{pmatrix}
    1&0&-3&|&0 \\ 0&2&-9&|&0 \\ 0&-1&5&|&1
\end{pmatrix}\]
Then \(R_3 \rightarrow 2R_3 + R_2\)
\[\begin{pmatrix}
    1&0&-3&|&0 \\ 0&2&-9&|&0 \\ 0&0&1&|&2
\end{pmatrix}\]
Hence \(i = 2\). 
\[2h-9i=0 \implies h =9\]
\[g-3i = 0 \implies g = 6\]
\[(g,h,i) = (6,9,2)\]
\subsection{}
\[A = \begin{pmatrix}
    1&0&-3 \\ 1&2&-12 \\ 1&-1&2
\end{pmatrix}\]
\[B=\begin{pmatrix}
    -8&3&6 \\ -14&5&9 \\ -3&1&2
\end{pmatrix}\]
\[AB = \begin{pmatrix}
    1&0&0 \\ 0&1&0 \\ 0&0&1
\end{pmatrix}\]
Hence AB is the identity matrix
\subsection{}
Each of the vectors are the standard basis vectors. Finding vectors that satisfy the conditions in parts a, b and c is the process to find the inverse of the matrix. Hence they have to multiply to the identity matrix by definition. 

\end{document}
